\[ %%%%%%%%%%%%%%%%%%%%%%%%%%%%%%%%%%%%%%%%%%%%%%%%%%%%%%%%%%%%%%%%%%%%%%%%%%%%%%%% \subsection{Long-Horizon Prediction Performance} \label{sec:long_horizon_results} %%%%%%%%%%%%%%%%%%%%%%%%%%%%%%%%%%%%%%%%%%%%%%%%%%%%%%%%%%%%%%%%%%%%%%%%%%%%%%%% The first evaluation focuses on the ability of PHYSGEN to perform accurate long-horizon prediction of nonlinear dynamical systems. Unlike one-step forecasting, long-horizon prediction requires a model to repeatedly apply its learned evolution operator: \begin{equation} \hat{X}_{t+k} = \Phi_\theta^k(X_t), \end{equation} where prediction errors can accumulate over time. The objective of this experiment is to determine whether physics-guided graph evolution reduces error propagation during autonomous rollout. %%%%%%%%%%%%%%%%%%%%%%%%%%%%%%%%%%%%%%%%%%%%%%%%%%%%%%%%%%%%%%%%%%%%%%%%%%%%%%%% \subsubsection{Evaluation Protocol} %%%%%%%%%%%%%%%%%%%%%%%%%%%%%%%%%%%%%%%%%%%%%%%%%%%%%%%%%%%%%%%%%%%%%%%%%%%%%%%% For each benchmark system, models are initialized from the same initial state: \begin{equation} X_0. \end{equation} The models then generate future trajectories without access to ground-truth states: \begin{equation} \hat{X}_{1}, \hat{X}_{2}, ... ,\hat{X}_{K}. \end{equation} The forecasting horizon is gradually increased: \begin{equation} K\in\{10,20,50,100\}. \end{equation} Performance is evaluated using the relative prediction error: \begin{equation} E_{rel}(k) = \frac{ ||X_{t+k}-\hat{X}_{t+k}||_2 }{ ||X_{t+k}||_2 }. \end{equation} %%%%%%%%%%%%%%%%%%%%%%%%%%%%%%%%%%%%%%%%%%%%%%%%%%%%%%%%%%%%%%%%%%%%%%%%%%%%%%%% \subsubsection{Overall Forecasting Accuracy} %%%%%%%%%%%%%%%%%%%%%%%%%%%%%%%%%%%%%%%%%%%%%%%%%%%%%%%%%%%%%%%%%%%%%%%%%%%%%%%% The results are summarized using the average long-horizon error: \begin{equation} LHE = \frac{1}{K} \sum_{k=1}^{K} E_{rel}(k). \end{equation} The expected comparison is: \begin{table}[h] \centering \caption{Long-horizon prediction performance comparison.} \label{tab:long_horizon_results} \begin{tabular}{lccc} \hline \textbf{Model} & \textbf{Short Horizon} & \textbf{Medium Horizon} & \textbf{Long Horizon} \\ \hline MLP & Low & Low & Poor \\ LSTM & Medium & Medium & Degraded \\ GNN & High & Medium & Degraded \\ GNS & High & High & Medium \\ PINN & Medium & Medium & Medium \\ FNO & High & High & Medium \\ DeepONet & High & High & Medium \\ PHYSGEN & High & High & High \\ \hline \end{tabular} \end{table} %%%%%%%%%%%%%%%%%%%%%%%%%%%%%%%%%%%%%%%%%%%%%%%%%%%%%%%%%%%%%%%%%%%%%%%%%%%%%%%% \subsubsection{Lorenz-63 Chaotic Dynamics} %%%%%%%%%%%%%%%%%%%%%%%%%%%%%%%%%%%%%%%%%%%%%%%%%%%%%%%%%%%%%%%%%%%%%%%%%%%%%%%% The Lorenz-63 experiment evaluates whether PHYSGEN can maintain stable prediction of chaotic trajectories. For chaotic systems, small approximation errors typically amplify rapidly: \begin{equation} D(t) \approx D_0e^{\lambda t}, \end{equation} where $\lambda$ represents the effective divergence rate. Models without explicit stability constraints are expected to exhibit faster trajectory divergence. PHYSGEN reduces divergence by combining: \begin{itemize} \item physics-guided message passing; \item latent state regularization; \item stability-aware optimization. \end{itemize} The resulting trajectory evolution is expected to preserve the qualitative structure of the attractor for longer horizons. %%%%%%%%%%%%%%%%%%%%%%%%%%%%%%%%%%%%%%%%%%%%%%%%%%%%%%%%%%%%%%%%%%%%%%%%%%%%%%%% \subsubsection{Reaction-Diffusion Prediction} %%%%%%%%%%%%%%%%%%%%%%%%%%%%%%%%%%%%%%%%%%%%%%%%%%%%%%%%%%%%%%%%%%%%%%%%%%%%%%%% For the Reaction-Diffusion benchmark, the evaluation focuses on spatial and temporal consistency. The predicted field: \begin{equation} \hat{u}(x,t) \end{equation} is compared with the reference solution: \begin{equation} u(x,t). \end{equation} The spatial prediction error is: \begin{equation} E_{space} = \frac{ ||u-\hat{u}||_2 }{ ||u||_2 }. \end{equation} The graph representation allows PHYSGEN to capture local interactions between neighboring spatial elements. %%%%%%%%%%%%%%%%%%%%%%%%%%%%%%%%%%%%%%%%%%%%%%%%%%%%%%%%%%%%%%%%%%%%%%%%%%%%%%%% \subsubsection{Navier--Stokes Forecasting} %%%%%%%%%%%%%%%%%%%%%%%%%%%%%%%%%%%%%%%%%%%%%%%%%%%%%%%%%%%%%%%%%%%%%%%%%%%%%%%% The Navier--Stokes experiment evaluates large-scale nonlinear fluid prediction. The model predicts: \begin{equation} \hat{X} = [\hat{u},\hat{v},\hat{p}], \end{equation} and is evaluated according to: \begin{itemize} \item velocity prediction error; \item pressure field accuracy; \item divergence constraint preservation. \end{itemize} The incompressibility condition: \begin{equation} \nabla\cdot u=0 \end{equation} is additionally monitored during rollout. %%%%%%%%%%%%%%%%%%%%%%%%%%%%%%%%%%%%%%%%%%%%%%%%%%%%%%%%%%%%%%%%%%%%%%%%%%%%%%%% \subsubsection{Error Growth Analysis} %%%%%%%%%%%%%%%%%%%%%%%%%%%%%%%%%%%%%%%%%%%%%%%%%%%%%%%%%%%%%%%%%%%%%%%%%%%%%%%% To analyze long-term behavior, the evolution of prediction error is visualized: \begin{equation} E(k) = ||X_{t+k}-\hat{X}_{t+k}||. \end{equation} The expected behavior is: \begin{itemize} \item data-driven models: rapid error growth; \item graph models: improved local prediction but increasing long-term drift; \item physics-informed models: better physical validity but possible optimization limitations; \item PHYSGEN: slower error accumulation due to stability regulation. \end{itemize} %%%%%%%%%%%%%%%%%%%%%%%%%%%%%%%%%%%%%%%%%%%%%%%%%%%%%%%%%%%%%%%%%%%%%%%%%%%%%%%% \subsubsection{Figure Representation} %%%%%%%%%%%%%%%%%%%%%%%%%%%%%%%%%%%%%%%%%%%%%%%%%%%%%%%%%%%%%%%%%%%%%%%%%%%%%%%% The main visualization is: \begin{figure}[h] \centering \includegraphics[width=0.9\linewidth]{figures/long_horizon_error.pdf} \caption{ Long-horizon prediction error growth comparison between PHYSGEN and baseline models. The x-axis represents prediction horizon $k$, and the y-axis represents relative prediction error. } \label{fig:long_horizon_error} \end{figure} %%%%%%%%%%%%%%%%%%%%%%%%%%%%%%%%%%%%%%%%%%%%%%%%%%%%%%%%%%%%%%%%%%%%%%%%%%%%%%%% \subsubsection{Discussion} %%%%%%%%%%%%%%%%%%%%%%%%%%%%%%%%%%%%%%%%%%%%%%%%%%%%%%%%%%%%%%%%%%%%%%%%%%%%%%%% The long-horizon prediction experiments demonstrate the importance of combining graph-based representation learning with explicit physical constraints. While existing approaches may achieve competitive short-term accuracy, their performance typically degrades as prediction horizons increase due to accumulated modeling errors. PHYSGEN addresses this limitation by learning a constrained dynamical operator that maintains both predictive accuracy and stable physical evolution. %%%%%%%%%%%%%%%%%%%%%%%%%%%%%%%%%%%%%%%%%%%%%%%%%%%%%%%%%%%%%%%%%%%%%%%%%%%%%%%% \subsubsection{Summary} %%%%%%%%%%%%%%%%%%%%%%%%%%%%%%%%%%%%%%%%%%%%%%%%%%%%%%%%%%%%%%%%%%%%%%%%%%%%%%%% The long-horizon evaluation provides the first validation of the central hypothesis of PHYSGEN. The results demonstrate that physics-guided graph evolution is a promising approach for stable prediction of nonlinear dynamical systems beyond short-term forecasting. \] 
